\documentclass[11pt]{article}
\usepackage[margin=1in]{geometry}
\usepackage[T1]{fontenc}
\usepackage[utf8]{inputenc}
\usepackage{amsmath,amssymb,amsthm}
\usepackage{enumitem}

\title{ICPC World Finals 2020\\K. Space Walls}
\author{}
\date{}

\begin{document}
\maketitle

\section*{Problem Summary}

We have a union of axis-aligned unit cubes. Each robot starts on an exposed unit face, facing along one of the
two tangent axis directions of that face, and then keeps moving ``straight'' over the surface. We must find the
earliest time when either
\begin{itemize}[leftmargin=*]
    \item two robots are on the same face at an integer time, or
    \item two robots move across the same pair of adjacent faces in opposite directions during one time step.
\end{itemize}

\section*{Key Observations}

\begin{itemize}[leftmargin=*]
    \item Let the current face normal be on axis $n$ and the current moving direction be on axis $d$.
    The third axis $a$ is neither $n$ nor $d$.
    During the whole motion, the coordinate on axis $a$ stays equal to its initial half-integer value.
    So every robot is confined to one fixed plane $a = c + \tfrac12$.
    \item Intersect the station with that plane. We obtain a union of unit squares on a 2D grid.
    Every exposed cube face whose normal is orthogonal to axis $a$ becomes one boundary unit edge of this 2D union.
    Moving straight on the 3D surface is exactly the same as walking along the boundary of this 2D union, one unit
    edge per time step.
    \item Therefore each robot moves on one directed boundary cycle of one 2D slice.
    After choosing a fixed orientation of every cycle, the position of a robot is just
    \[
        s + \delta t \pmod L,
    \]
    where $L$ is the cycle length, $s$ is the starting offset on the cycle, and $\delta \in \{+1,-1\}$ tells whether
    the robot follows the chosen orientation or the opposite one.
    \item If two robots use the same invariant axis and the same slice, then they can collide only if they are on the
    same boundary cycle.
    If their invariant axes are different, then they can never swap locations; they can only meet on a common face.
\end{itemize}

\section*{Algorithm}

\subsection*{1. Build only the slices that are actually needed}

The invariant axis and slice coordinate of a robot are determined by its input face normal and initial moving
direction. There are at most $k$ distinct slices, so we build only those.

For one fixed slice:
\begin{enumerate}[leftmargin=*]
    \item Take all station boxes intersecting that slice.
    \item Project them onto the slice plane.
    \item Coordinate-compress the rectangle borders on the two slice axes.
    \item Mark the occupied compressed cells.
\end{enumerate}

\subsection*{2. Extract directed boundary cycles}

For every occupied compressed cell, each side adjacent to empty space contributes one directed primitive boundary edge.
The direction is chosen so that the occupied cell lies on the \emph{left} side of the edge.

At the end of a directed edge, the next edge on the same contour is found by the usual left-hand rule:
try \emph{left}, then \emph{straight}, then \emph{right}, then \emph{back}.
This decomposes all primitive edges into disjoint directed cycles.

Along each cycle we merge consecutive collinear primitive edges into one segment and store:
\begin{itemize}[leftmargin=*]
    \item its direction,
    \item the exposed face normal represented by this segment,
    \item its coordinate interval,
    \item its prefix length on the cycle.
\end{itemize}

\subsection*{3. Map every starting robot state to a cycle position}

The starting face is one unit sub-edge of exactly one segment in its slice.
Hence we can determine:
\begin{itemize}[leftmargin=*]
    \item which cycle contains the robot,
    \item the cycle length $L$,
    \item the starting index $s$ of its current face on that cycle,
    \item whether its motion agrees with the stored cycle orientation ($\delta = +1$) or not ($\delta = -1$).
\end{itemize}

\subsection*{4. Check every pair of robots}

\paragraph{Same invariant axis and same slice.}
They can interact only on the same cycle.

If the cycle is the same, then the face occupied by robot $i$ at time $t$ is
\[
    p_i(t) \equiv s_i + \delta_i t \pmod L.
\]

\begin{itemize}[leftmargin=*]
    \item Same-face collision: solve
    \[
        s_1 + \delta_1 t \equiv s_2 + \delta_2 t \pmod L.
    \]
    \item Swap collision: this is possible only when $\delta_1 = -\delta_2$.
    Then robot $2$ must be exactly one step ahead of robot $1$ along robot $1$'s moving direction:
    \[
        p_2(t) - p_1(t) \equiv \delta_1 \pmod L.
    \]
\end{itemize}

Both are linear congruences in one variable.

\paragraph{Different invariant axes.}
Assume robot $A$ uses axis $a$ and robot $B$ uses axis $b$, and let $c$ be the third axis.
They can meet only on faces orthogonal to axis $c$.

For a fixed robot, a cycle segment representing such faces runs along the other robot's invariant axis, so after fixing
the other robot's slice coordinate, that whole segment contributes at most one candidate face.

For robot $A$, scan all its cycle segments and collect all candidate common faces into a hash map:
\[
    (\text{plane coordinate on axis } c,\ \text{sign of the normal}) \mapsto \text{time residue modulo } L_A.
\]
Do the same scan for robot $B$; whenever the same face appears in both scans, solve the two congruences
\[
    t \equiv r_A \pmod{L_A}, \qquad t \equiv r_B \pmod{L_B}
\]
with CRT and keep the smallest solution.

\section*{Correctness Proof}

We prove that the algorithm returns the earliest collision time.

\paragraph{Lemma 1.}
For any robot, the coordinate on exactly one axis stays constant during the whole motion.

\paragraph{Proof.}
At every step, the current face normal axis and the current moving axis are two different coordinate axes.
The robot always moves over an edge shared by those two directions, so the third axis is not involved in the motion.
When the robot continues onto the next face, both the old face center and the new face center still have the same
half-integer coordinate on that third axis. Thus that coordinate is invariant forever. \qed

\paragraph{Lemma 2.}
Inside one fixed slice plane, the robot motion is exactly a walk along a boundary cycle of the 2D cross-section, moving
one boundary unit edge per time step.

\paragraph{Proof.}
Consider the intersection of the station with the invariant slice plane from Lemma~1.
Every cube cut by the plane contributes one unit square, so the cross-section is a union of grid squares.
An exposed 3D face orthogonal to the slice plane becomes a boundary unit edge of this union.

Now inspect one robot step locally near the crossed cube edge.
In the slice plane, this step is precisely the transition from one boundary unit edge to the next boundary unit edge of
the cross-section contour.
The ``cannot slip between cubes'' rule is exactly the statement that the walk follows the geometric boundary of the
occupied region, not a nonexistent gap through it.
Therefore the robot walks along one boundary cycle, one unit edge per time step. \qed

\paragraph{Lemma 3.}
After fixing an orientation of a boundary cycle, a robot on that cycle occupies position
$s + \delta t \pmod L$ at time $t$ for some constants $s$, $\delta \in \{+1,-1\}$ and cycle length $L$.

\paragraph{Proof.}
By Lemma~2 the robot advances by exactly one boundary unit edge per time step on one cycle of length $L$.
If the chosen cycle orientation matches the robot's actual traversal direction, the offset increases by $1$ each step;
otherwise it decreases by $1$ each step. Hence the formula follows. \qed

\paragraph{Lemma 4.}
For two robots on the same slice, the algorithm finds exactly all same-face and swap collisions between them.

\paragraph{Proof.}
Different boundary cycles of the same slice correspond to disjoint sets of exposed faces, so only robots on the same
cycle can interact.
For robots on one common cycle, Lemma~3 gives explicit formulas for both positions as functions of time.

A same-face collision happens exactly when the two formulas are equal modulo $L$, which is the first congruence solved
by the algorithm.

A swap collision on one cycle happens exactly when, at time $t$, robot $2$ occupies the unique neighboring cycle edge
that robot $1$ will enter at time $t+1$.
That is exactly the second congruence used by the algorithm.
So every and only valid same-slice collision is detected. \qed

\paragraph{Lemma 5.}
For two robots with different invariant axes, the algorithm finds exactly all collisions between them.

\paragraph{Proof.}
Such robots move in different slice planes, so a swap is impossible: one step between adjacent faces is always contained
in a single slice plane.
Therefore only same-face collisions matter.

Let the invariant axes be $a$ and $b$, and let $c$ be the third axis.
Any common face must be orthogonal to axis $c$, because that is the only face type lying in both slice planes.
For each robot, every relevant cycle segment yields at most one such face after fixing the other robot's slice
coordinate. The algorithm enumerates precisely those faces and computes, via Lemma~3, the unique time residue modulo the
robot's period when that face is occupied.

Thus a collision on a particular common face exists exactly when the two residues are compatible, which is exactly what
the CRT check tests. Therefore the algorithm finds every and only valid cross-slice collision. \qed

\paragraph{Theorem.}
The algorithm outputs the earliest collision time, or \texttt{ok} if no collision ever occurs.

\paragraph{Proof.}
By Lemmas~4 and~5, for every pair of robots the algorithm enumerates exactly all collision times involving that pair and
keeps the smallest one. Any collision among several robots contains at least one colliding pair, so the minimum over all
pairs is the global earliest collision time. If no pair yields a collision time, then no collision ever occurs, so
\texttt{ok} is correct. \qed

\section*{Complexity Analysis}

Let $S$ be the number of distinct slices used by the robots; clearly $S \le k$.

For one slice, there are at most $2n$ compressed coordinates on each slice axis, hence $O(n^2)$ compressed cells and
$O(n^2)$ primitive boundary edges.
Building one slice with the straightforward marking used in the implementation takes $O(n^3)$ time in the worst case
and $O(n^2)$ memory.

For one pair of robots, checking collisions scans only the segments on their own cycles, which is $O(n^2)$ in the worst
case.

Therefore the total complexity is
\[
    O(Sn^3 + k^2 n^2)
\]
time and
\[
    O(n^2)
\]
memory per slice.
With $n,k \le 100$, this easily fits.

\section*{Implementation Notes}

\begin{itemize}[leftmargin=*]
    \item The cycle lengths can be large, and the CRT may combine two periods, so use 64-bit integers.
    \item Boundary edges are directed so that occupied area stays on the left; this makes the successor rule uniform.
    \item In the cross-slice case, the key of a common face is just
    \[
        (\text{coordinate of the plane orthogonal to the third axis},\ \text{sign of the normal}),
    \]
    because the two slice coordinates are already fixed by the robot pair.
\end{itemize}

\end{document}
